\documentclass{article}
\usepackage[T2A]{fontenc}
\usepackage[english, russian]{babel}
\usepackage{mathtools, amssymb, amsfonts, amsthm}
\usepackage{tikz-cd}

\newtheorem*{theorem}{Теорема}
\newtheorem*{corollary}{Следствие}
\newtheorem*{lemma}{Лемма}
\newtheorem*{statement}{Утверждение}
\newtheorem*{proposition}{Предложение}
\theoremstyle{definition}
\newtheorem*{definition}{Определение}
\theoremstyle{remark}
\newtheorem*{reminder}{Напоминание}
\newtheorem*{remark}{Замечание}
\newtheorem*{example}{Пример}
\newtheorem{problem}{Задача}

\DeclareMathOperator{\chr}{chr}
\DeclareMathOperator{\disc}{disc}

\newcommand{\llangle}{\langle \hspace{-.5mm} \langle}
\newcommand{\rrangle}{\rangle \hspace{-.5mm} \rangle}

\title{Квадратичные формы над полями}
\author{}
\date{}

\begin{document}
\maketitle

Список литературы:

\begin{enumerate}
    \item \foreignlanguage{english}{Lam. Introduction to quadratic forms over fields.}
    \item \foreignlanguage{english}{Scharlau. Quadratic and Hermitian forms}
    \item \foreignlanguage{english}{Elman, Karpenko, Merkurjev. The algebraic and geometric theory of quadratic forms.}
    \item \foreignlanguage{english}{Serre. A course in arithmetic.}
    \item Шафаревич. Теория чисел.
    \item Касселс. Рациональные квадратичные формы.
\end{enumerate}

\begin{definition}
$F$ --- поле, $\chr F \neq 2$, $V/ F$ --- конечномерное векторное пространство. Отображение $q: V \longrightarrow F$ называется квадратичной формой, если 
\begin{gather*}
q(av) = a^2 q(v) \\
B(u, v) = \frac 1 2 (q(u + v) - q(u) - q(v))\text{ --- билинейная форма}
\end{gather*}
\end{definition}

Нетрудно заметить, что существуеут взаимнооднозначное соответсвие между квадратичными формами и билинейными симметричными формами.

Пусть $e_1, e_2, \ldots, e_n$ --- базис $V$. Тогда 

$$q(x_1e_1 + x_2 e_2 + \ldots + x_n e_n) = \sum B(e_i, e_j)x_ix_j$$

Матрица $(B(e_i, e_j))$ называется матрицой квадратичной формы $q$ в базисе $\{e_i\}$. Нетрудно видеть, что при замене базиса $A \rightsquigarrow S^T A S$, так что $\det q \in F^* / F^{*2}$ не зависит от выбора базиса.

Квадратичная форма задаёт линейное отображение в двойственное пространство.

$$f: V \rightarrow V^*: v \rightarrow (u \mapsto B(u, v))$$

\begin{proposition}
    $f$ --- изоморфизм тогда и только тогда, когда $B$ невырождена.
\end{proposition}

\begin{definition}
    $q$ называется изотропной, если существует ненулевой вектор $v$ такой, что $q(v) = 0$. В противном случае $q$ называется анизотропной.
\end{definition}

\begin{definition}
    Пусть $(V_1, q_1)$, $(V_2, q_2)$ --- пространства с квадратичными формами. Тогда на $V_1 \oplus V_2$ определена квадратичная форма $q_1 \bot q_2(v_1, v_2) := q_1(v_1) + q_2(v_2)$.
\end{definition}

\begin{definition}
    Аналогично на $V_1 \otimes V_2$ определена квадратчиная форма $q_1 \otimes q_2(v_1 \otimes v_2) = q_1(v_1) q_2(v_2)$.
\end{definition}

\begin{definition}
    Формы $q_1$ и $q_2$ называются изоморфными если существует изоморфизм $f: V_1 \rightarrow V_2$ такой, что $q_1(v_1) = q_2(f(v_1))$.
\end{definition}

\begin{theorem}[Теорема о сокращении]
    Если $f_1 \bot g \simeq f_2 \bot g$, то и $f_1 \simeq f_2$.
\end{theorem}

\begin{theorem}
    Пусть $V_0 \subset V$ --- подпространство, $q_0 = q|V$ --- подформа формы $q$. Если $q_0$ невырождена, то $q \simeq q_0 \bot q_1$.
\end{theorem}
\begin{proof}
    $V_0^\bot = \{v \in V | B_q(v, v_0) = 0 \forall v_0 \in V_0\}$
    Тогда $V_0^\bot \cap V_0 = 0$. Ну действительно, пусть $u$ лежит в пересечении, тогда $B(u, v_0) = 0$ для любого $v_0 \in v$, но тогда образ $u$ в $V_0^*$ равен нулю, значит, т. к. $q_0$ невырождена, $u = 0$.
    Остаётся заметить, что существует точная последовательность
    \begin{equation*}
    \begin{tikzcd}
    0 \rar & V_0^\bot \rar & V \rar & V_0^*
    \end{tikzcd}
    \end{equation*}
    Значит $\dim V = \dim V_0^* + \dim V_0^\bot = \dim V_0 + \dim V_0^\bot$, а значит $V = V_0 \oplus V_0^\bot$.
\end{proof}

\begin{definition}
    Кольцо Витта $W(F)$. Пусть $f_1$, $f_2$ --- невырожденные квадратичные формы, если существуют $m, n \geq 0$ такие, что $f_1 \bot mH \simeq f_2 \bot nH$, где $H = \langle 1, -1 \rangle$ --- гиперболическая плоскость.
    $$[f] + [g] = [f \bot g]$$
    $$[f] \otimes [g] = [f \otimes g]$$
\end{definition}

$W(F) \supset I(F)$ --- идеал классов чётномерных форм.

$W(F) \supset I(F) \supset I^2(F) \supset \ldots$

Вопросы:
\begin{enumerate}
    \item $\cap I^k(F) = 0$
    \item $I^n(F) / I^{n + 1} = \,?$
\end{enumerate}

$$I(F) \rightarrow F^* / F^{*2}$$
$$\disc f := (-1)^{\frac{n(n - 1)}{2}}$$

\begin{statement}
    $\ker \disc = I^2(F)$
\end{statement}

\begin{definition}
    Пусть $\varphi$ --- форма, $D(\varphi) \subset F^*$ --- множество ненулевых значений $\varphi$. $G(\varphi) = \{a \in F^*| \varphi \simeq a \varphi\}$. $G(\varphi)$, очевидно, является подгруппой. Если $1 \in D(\varphi)$, $G(\varphi) \subset D(\varphi)$
\end{definition}

\begin{definition}
    Формой Пфистера называется форма $\llangle a_1, a_2, \ldots, a_n \rrangle := \langle 1, -a_1 \rangle \otimes \langle 1, -a_2 \rangle \otimes \langle 1, -a_3 \rangle \otimes \ldots \otimes \langle 1, -a_n \rangle$.
\end{definition}

\begin{definition}
    $u(F)$ --- максимальная размерность анизотропной квадратчиной формы над $F$.
\end{definition}

\end{document}
